\chapter{Quantum Field Theory}
\label{chap:qft}

Quantum field theory (QFT) provides the modern theoretical framework for
describing elementary particles and their interactions. It underlies the
Standard Model of particle physics and plays a central role not only in
high-energy physics, but also in condensed matter theory and the physics of
many-body systems. In the context of dense matter and compact stars, QFT offers
a natural language for describing interacting particles at high densities and
for connecting microscopic physics to macroscopic quantities
\cite{PeskinSchroeder1995,Weinberg1995}.

The need for QFT arises from the limitations of
non-relativistic quantum mechanics. While quantum mechanics successfully
describes a wide range of phenomena at atomic and molecular scales, it is not
compatible with special relativity in a fundamental way. In particular, the
Schrödinger equation treats time and space asymmetrically and does not enforce
relativistic causality. When relativistic effects become important, the notion
of a fixed number of particles breaks down, and attempts to combine quantum
mechanics with special relativity at the level of single-particle wavefunctions
lead to conceptual and physical inconsistencies
\cite{BjorkenDrell1964}.

These issues can be illustrated with simple thought experiments. An
example is the relativistic boosting of a particle confined
to a box. When such a system is analyzed in a highly boosted reference frame,
quantum fluctuations in momentum lead to large fluctuations in energy that
cannot be consistently interpreted within ordinary quantum mechanics. In
particular, the appearance of energy fluctuations associated with the vacuum
indicates the need for a framework in which particle creation and annihilation
are allowed and the vacuum itself carries dynamical degrees of freedom. This
example provides a concrete and physically transparent motivation for QFT \cite{Vutha_2013}.

QFT resolves these problems by promoting particles to
excitations of underlying quantum fields that extend throughout space and time.
In this framework, each particle species is associated with a corresponding
field, and particles appear as localized excitations of that field. The field
theoretic description naturally incorporates special relativity and provides a
consistent treatment of processes in which particle number is not conserved.
It also offers a fundamental explanation for the existence of identical
particles: particles of the same type are identical because they arise from the
same underlying quantum field \cite{Weinberg1995}.

One of the most successful realizations of QFT is quantum
electrodynamics (QED), the relativistic field theory of electrons and photons.
QED is among the most precisely tested theories in physics, with experimental
measurements agreeing with theoretical predictions to remarkable accuracy. Its
success provides strong evidence that quantum fields, rather than particles
alone, constitute the fundamental degrees of freedom of nature
\cite{PeskinSchroeder1995}. More generally, QFT provides a
unified framework in which quantum mechanics and special relativity are combined
in a consistent and predictive manner.

For the study of dense matter, QFT is particularly useful
because it allows interactions, symmetries, and conservation laws to be treated
systematically. Thermodynamic quantities such as pressure and energy density can
be derived directly from an underlying Lagrangian, making it possible to
construct equations of state from microscopic models. In the following
sections, the basic concepts of QFT relevant for this work will
be introduced, with an emphasis on those needed to model strongly interacting
matter at finite density and to prepare for the QCD-inspired descriptions used
later in the thesis.

\section{Bridging quantum mechanics and field theory}
\label{sec:bridging_qm_qft}

In this section we briefly introduce the conceptual and mathematical framework
that connects ordinary quantum mechanics to QFT. The emphasis
is on the path-integral formulation, which generalizes naturally from systems
with finitely many degrees of freedom to relativistic fields and provides a
convenient starting point for the description of interacting matter.

\subsection*{Action principle and path integrals in quantum mechanics}

In classical mechanics, the dynamics of a system with generalized coordinate
$q(t)$ is determined by the action
\begin{equation}
S[q] = \int_{t_i}^{t_f} \! dt \, L\bigl(q(t), \dot{q}(t)\bigr),
\end{equation}
where $L$ is the Lagrangian of the system. The classical equations of motion
follow from the principle of stationary action, $\delta S = 0$.

In quantum mechanics, an alternative formulation of dynamics is provided by the
path-integral approach. Instead of describing time evolution in terms of wave
functions and operators, transition amplitudes between initial and final states
are expressed as sums over all possible trajectories connecting them. Formally,
the transition amplitude can be written as
\begin{equation}
\langle q_f, t_f | q_i, t_i \rangle
= \int \mathcal{D}q(t)\,
\exp\!\left(\frac{i}{\hbar} S[q]\right),
\end{equation}
where the functional integral runs over all paths $q(t)$ satisfying the
boundary conditions $q(t_i)=q_i$ and $q(t_f)=q_f$
\cite{FeynmanHibbs1965,Schulman1981}.

Physical observables in this formulation are expressed in terms of correlation
functions, which measure how the value of a dynamical variable at one time is
related to its value at another time. Correlation functions therefore encode
information about the system’s dynamics. In quantum mechanics, the dynamical 
variables are represented by operators, and expectation values are taken with
respect to a given quantum state. As an example, the time-ordered two-point function
\(\langle T\, q(t)\, q(t') \rangle\) describes how fluctuations of the quantum
degree of freedom \(q\) at time \(t'\) are correlated with fluctuations at time
\(t\). Here \(T\) denotes the time-ordering operator, which arranges operators
according to their time arguments and ensures a causal interpretation of the
correlation function. This two-point function therefore plays the role of a
propagator, characterizing how disturbances of the system evolve in time, and
this viewpoint generalizes directly to QFT, where the dynamical
variables are fields defined throughout spacetime.

\subsection*{Generating functionals}

A convenient way to handle correlation functions is to introduce an external
source that couples to the system. In quantum mechanics, this is done by adding
a term of the form $J(t)\,q(t)$ to the action. The source $J(t)$ is a classical,
prescribed function that does not describe a physical force, but is introduced
as a mathematical tool. It allows us to probe how the system responds to small
disturbances and to extract correlation functions in a systematic way.

With the source included, one defines the generating functional
\begin{equation}
Z[J] = \int \mathcal{D}q(t)\,
\exp\!\left[\frac{i}{\hbar}
\left(S[q] + \int dt\, J(t)\, q(t)\right)\right].
\end{equation}
By taking functional derivatives of $Z[J]$ with respect to the source and then
setting $J=0$, one generates time-ordered correlation functions of the dynamical
variable $q(t)$,
\begin{equation}
\langle T\, q(t_1)\cdots q(t_n)\rangle
= \left. \frac{1}{i^n}
\frac{\delta^n Z[J]}{\delta J(t_1)\cdots \delta J(t_n)}
\right|_{J=0}.
\end{equation}
In this sense, the generating functional acts as a compact object from which all
correlation functions of the theory can be obtained. This idea extends directly
to QFT, where the dynamical variables are fields and the source
becomes a function of spacetime.


\subsection*{From particles to fields}

The transition from quantum mechanics to QFT amounts to
replacing a finite set of dynamical variables by fields defined at every point
in spacetime. A real scalar field $\phi(x)$ assigns a single real number to each
spacetime point $x=(t,\mathbf{x})$. The dynamics of the field are specified by
a Lagrangian density $\mathcal{L}(\phi,\partial_\mu\phi)$, and the action takes
the form
\begin{equation}
S[\phi] = \int d^4x \, \mathcal{L}(\phi,\partial_\mu\phi).
\end{equation}
The classical equations of motion follow from the condition $\delta S=0$ and are
given by the Euler–Lagrange equations for fields.

In the path-integral formulation of QFT, the generating
functional is defined as
\begin{equation}
Z[J] = \int \mathcal{D}\phi \,
\exp\!\left[i\left(S[\phi] + \int d^4x\, J(x)\,\phi(x)\right)\right],
\end{equation}
where $J(x)$ is a spacetime-dependent source. As in quantum mechanics, all
time-ordered correlation functions of the field can be obtained by functional
differentiation with respect to $J(x)$. In particular, the two-point function
plays the role of the propagator and describes how field fluctuations propagate
through spacetime.

The path-integral formulation provides a unified and flexible framework for
QFT. It treats space and time on an equal footing, makes the
role of symmetries transparent, and generalizes naturally to systems with many
interacting degrees of freedom. For these reasons, it will serve as the primary
language for the scalar field theories discussed in the following sections and
for the effective models of dense matter introduced later in this thesis
\cite{PeskinSchroeder1995,Weinberg1995}.

% \section{Free scalar field}
% \label{sec:free_scalar_field}

% We now turn to the simplest example of a relativistic quantum field theory: a
% real scalar field. Despite its simplicity, this theory already illustrates
% many of the structural features common to more realistic models, including the
% role of symmetries, conserved currents, and propagators. It also provides a
% natural starting point for introducing interactions in a controlled way.

% \subsection*{Classical scalar field}

% A real scalar field $\phi(x)$ is a real-valued function defined on spacetime,
% where $x=(t,\mathbf{x})$. The dynamics of a free relativistic scalar field are
% specified by the Lagrangian density
% \begin{equation}
% \mathcal{L}_0
% = \frac{1}{2}\partial_\mu \phi\, \partial^\mu \phi
% - \frac{1}{2} m^2 \phi^2 ,
% \end{equation}
% where $m$ is a mass parameter and natural units $\hbar=c=1$ are used. The action
% is obtained by integrating the Lagrangian density over spacetime,
% \begin{equation}
% S_0[\phi] = \int d^4x \, \mathcal{L}_0 .
% \end{equation}

% Varying the action with respect to the field yields the Euler--Lagrange equation
% \begin{equation}
% \frac{\partial \mathcal{L}_0}{\partial \phi}
% - \partial_\mu
% \frac{\partial \mathcal{L}_0}{\partial(\partial_\mu\phi)} = 0 ,
% \end{equation}
% which leads to the Klein--Gordon equation
% \begin{equation}
% (\Box + m^2)\phi(x) = 0 ,
% \end{equation}
% where $\Box = \partial_\mu\partial^\mu$ is the d'Alembert operator. This equation
% describes a relativistic scalar particle of mass $m$.

% \subsection*{Symmetries and Noether's theorem}

% The action $S_0[\phi]$ possesses several continuous symmetries. Of particular
% importance are spacetime translations,
% \begin{equation}
% x^\mu \rightarrow x^\mu + a^\mu ,
% \end{equation}
% under which the field transforms as $\phi(x)\rightarrow \phi(x-a)$. According
% to Noether's theorem, every continuous symmetry of the action gives rise to a
% conserved current \cite{Noether1918,Weinberg1995}. For spacetime translations,
% the associated conserved quantity is the energy--momentum tensor,
% \begin{equation}
% T^{\mu\nu}
% = \partial^\mu \phi\, \partial^\nu \phi
% - g^{\mu\nu}\mathcal{L}_0 ,
% \end{equation}
% which satisfies the conservation law $\partial_\mu T^{\mu\nu}=0$ when the
% equations of motion are fulfilled. The conserved charges obtained by integrating
% $T^{0\nu}$ over space correspond to the total energy and momentum of the field.

% The connection between symmetries and conservation laws provided by Noether's
% theorem plays a central role in quantum field theory. In later chapters, it
% will be used to identify conserved charges associated with internal symmetries
% and to motivate the introduction of chemical potentials in field-theoretic
% descriptions of dense matter.

% \subsection*{Free scalar quantum field and generating functional}

% The quantum theory of the free scalar field is defined through the path
% integral. Introducing a spacetime-dependent source $J(x)$, the generating
% functional is given by
% \begin{equation}
% Z_0[J] =
% \int \mathcal{D}\phi\,
% \exp\!\left[i\left(
% S_0[\phi] + \int d^4x\, J(x)\phi(x)
% \right)\right].
% \end{equation}
% Because the action is quadratic in the field, the functional integral is
% Gaussian and can be evaluated exactly. While the explicit calculation will not
% be carried out here, its structure implies that all correlation functions can be
% expressed in terms of the two-point function.

% The time-ordered two-point correlation function,
% \begin{equation}
% D_F(x-y)
% = \langle T\,\phi(x)\phi(y)\rangle ,
% \end{equation}
% is known as the Feynman propagator. It describes the propagation of field
% fluctuations between spacetime points $x$ and $y$. Formally, the propagator is
% the Green's function of the Klein--Gordon operator,
% \begin{equation}
% (\Box_x + m^2) D_F(x-y) = -i\,\delta^{(4)}(x-y) .
% \end{equation}
% All higher-order correlation functions of the free theory can be obtained by
% functional differentiation of $Z_0[J]$ with respect to the source $J(x)$.

% The free scalar field thus provides a complete and exactly solvable example of a
% relativistic quantum field theory. Its simplicity makes it an ideal laboratory
% for introducing interactions and symmetry-breaking mechanisms, which will be
% the focus of the next section.

% \section{Interacting scalar fields: the \texorpdfstring{$\phi^4$}{phi4} theory}
% \label{sec:phi4_theory}

% Having established the basic framework for free scalar fields, we now introduce
% interactions in the simplest possible setting. The interacting scalar field
% theory with a quartic self-interaction, commonly referred to as $\phi^4$ theory,
% serves as a prototypical example for studying the effects of interactions and
% spontaneous symmetry breaking in quantum field theory. Although highly
% idealized, this model captures several structural features that also appear in
% more realistic theories of strongly interacting matter.

% \subsection*{Lagrangian and symmetries}

% The Lagrangian density of the interacting scalar field is given by
% \begin{equation}
% \mathcal{L}
% = \frac{1}{2}\partial_\mu \phi\,\partial^\mu \phi
% - \frac{1}{2} m^2 \phi^2
% - \frac{\lambda}{4!}\,\phi^4 ,
% \end{equation}
% where $\lambda$ is a real coupling constant. For $\lambda>0$, the interaction
% term ensures that the energy of the system is bounded from below.

% The theory is invariant under the discrete transformation
% \begin{equation}
% \phi(x) \rightarrow -\phi(x) ,
% \end{equation}
% which corresponds to a $\mathbb{Z}_2$ symmetry of the action. This symmetry plays
% a central role in the discussion of spontaneous symmetry breaking.

% \subsection*{Classical potential and vacuum structure}

% For static and spatially uniform field configurations, the dynamics of the
% theory are governed by the classical potential
% \begin{equation}
% V(\phi)
% = \frac{1}{2} m^2 \phi^2
% + \frac{\lambda}{4!}\,\phi^4 .
% \end{equation}
% The vacuum of the theory corresponds to the field configuration that minimizes
% this potential.

% If $m^2>0$, the potential has a single minimum at $\phi=0$, and the $\mathbb{Z}_2$
% symmetry is realized trivially. In contrast, if $m^2<0$, the potential takes the
% form of a double well with two degenerate minima located at
% \begin{equation}
% \phi = \pm v, \qquad
% v = \sqrt{\frac{6|m^2|}{\lambda}} .
% \end{equation}
% In this case, the ground state of the theory is not invariant under the
% $\mathbb{Z}_2$ symmetry of the Lagrangian.

% \subsection*{Spontaneous symmetry breaking}

% When the system selects one of the two degenerate minima as its vacuum state,
% the discrete symmetry $\phi\rightarrow-\phi$ is said to be spontaneously broken.
% Although the Lagrangian remains symmetric, the vacuum does not share this
% symmetry. This phenomenon is referred to as spontaneous symmetry breaking and
% occurs whenever the ground state of a system has a lower symmetry than the
% underlying equations of motion.

% To study fluctuations around a chosen vacuum, one writes the field as
% \begin{equation}
% \phi(x) = v + \eta(x) ,
% \end{equation}
% where $\eta(x)$ describes small fluctuations about the vacuum expectation value
% $v$. Substituting this expression into the Lagrangian and expanding in powers of
% $\eta$ yields a new Lagrangian for the fluctuation field. To quadratic order,
% one finds a mass term for $\eta$ with
% \begin{equation}
% m_\eta^2 = 2|m^2| ,
% \end{equation}
% indicating that the excitations around the vacuum correspond to massive scalar
% particles.

% This simple example illustrates how interactions can dynamically generate mass
% scales and nontrivial vacuum structure, even when the original Lagrangian
% contains only a small number of parameters.

% \subsection*{Interpretation and relevance}

% Although $\phi^4$ theory is not intended as a realistic model of particle
% physics, it plays an important conceptual role. It provides a minimal setting
% in which interactions, symmetry, and vacuum structure can be analyzed
% explicitly. In particular, it demonstrates how spontaneous symmetry breaking
% arises from the interplay between the form of the potential and the symmetries
% of the action.

% These ideas extend far beyond scalar field theory. In more realistic models,
% such as those used to describe strongly interacting matter, scalar fields often
% serve as effective degrees of freedom associated with order parameters of
% symmetry-breaking transitions. The $\phi^4$ theory thus serves as a useful
% prototype for understanding the structure of effective field-theoretic models
% employed later in this thesis, including those inspired by quantum
% chromodynamics.
